%------------------------------------------------------------------------------%
% \chapter{Plano Cartesiano e Equações}
%------------------------------------------------------------------------------%

%------------------------------------------------------------------------------%
\begin{exercicios}{Posição Relativa de Retas}
    
    \begin{exercicio}
        Se a reta que passa pelos pontos $(1,a)$ e $(4,-2)$ é paralela a
        reta que passa pelos pontos $(2,8)$ e $(-7,a+4)$ qual o valor de $a$?
    \end{exercicio}
    
    \begin{exercicio}
        Determinar o ponto onde as retas $3y+3x-2=0$ e $3y-x=1$ se cruzam.
    \end{exercicio}
    
    \begin{exercicio}
        Encontre uma equação para a reta que passa pelo ponto $(2, 3)$ e é  paralela à reta com equação $3x + 4y - 8 = 0$. 		
    \end{exercicio}
    
    \begin{exercicio}
        Encontre uma equação para a reta que passa pelo ponto $(-1, 3)$ e é paralela à reta que passa pelos pontos $(-3, 4)$ e $(2, 1)$.	
    \end{exercicio}

    \begin{exercicio}
        Encontre uma equação para a reta que passa pelo ponto $(-2,-4)$ e é perpendicular à reta com equação $2x-3y-24=0$. 
    \end{exercicio}

    \begin{exercicio}
        Encontre uma equação para a reta que passa pelo ponto $(-2,-4)$ e é perpendicular à reta com equação $2x-3y-24=0$.
    \end{exercicio}
    
    \begin{exercicio}
        Construa a equação da reta $s$ que passa pelo ponto $(1,-1)$ e é perpendicular a reta $r$, definida pela equação $2y=x+4$.
    \end{exercicio}

    \begin{exercicio}
        Determine a interseção entre as retas $x+2y=3$ e $2x+3y=5$.
    \end{exercicio}
    
    \begin{exercicio}
        Mostrar que as retas $2x+3y-1=0$, $x+y=0$ e $3x+4y-1=0$ se encontram no mesmo ponto.
    \end{exercicio}
    
\end{exercicios}

%------------------------------------------------------------------------------%
